
\section{Conclusions}
The central conclusion of this work is that treating time as a three-dimensional vector subspace remains compatible with established relativistic physics while providing a geometric origin for structures conventionally introduced as independent elements of quantum mechanics and quantum field theory. Within this architecture, the arrow of time, complex phase, intrinsic spin, wave--particle complementarity, gauge covariance, Abelian and non-Abelian gauge structures, mass, scalar Yukawa interactions, charged-pseudoscalar isospin exchange, and anomalous magnetic-moment interactions arise from distinct aspects of the same temporal geometry.

Taken together, these results answer the guiding question of this work affirmatively and thereby fulfill its stated scope. More broadly, they lay the foundations for a consistent theory of quantum physics in which Einsteinian relativity and quantum structures arise together from a common six-dimensional geometry, rather than quantum dynamics being independently postulated and subsequently constrained to conform to a pre-existing relativistic spacetime. This common origin distinguishes the present unification from existing approaches, while the framework's economy and broad results offer a promising route to general unification.
